Firms reallocate capacity as market opportunities change, altering competitors'
profitability and the cash flows investors anticipate.
We extend the multi-contest framework of
\citet{xu_zenou_zhou_conflict_2022} to explain how competition and adjustment
costs shape firms' changing positions and their exposure to free cash flow
disturbances.
The model establishes equilibrium properties and comparative statics, linking
gradual operating adjustment to financial responses when new economic
opportunities arrive.

In a restricted two-opportunity economy, firms with greater initial exposure
experience larger return responses to positive exogenous free cash flow shocks.
As other firms expand toward that opportunity, their news distributions
approach the initial distributions of those competitors.
Under common reporting intensity and the stated timing conditions, this
mechanism implies a nonnegative expected local association between financial
returns and movement toward competitors in information space.

We develop a novel sandwich estimator that relates changes in Wasserstein
distances between
corporate-news distributions to excess returns weighted by article timing.
Multi-agent simulations show how limited coverage obscures small operating
adjustments and how changing coverage can generate similar observed movements.
In an application using Nasdaq-listed firms our empirical study finds a
small positive association across alternative specifications.
Together, these results give an economic foundation for studying firms'
changing economic footprints through financial news and clarifies how competition
and information production shape firms' investment decisions.
